% Propietario conservado: /Users/ruben/Documents/ChatGPT/jueces y controles/REVISION_ARGUMENTAL_APERTURAS_CIERRES_20260915/delivery/ARTICLE_V_ES_ARGUMENTAL_20260915/manuscrito/sections/v_coordenadas_angulares.tex
% RESULTADO_RECUPERADO desde II REV09. Véase MAPA_PROCEDENCIA.json.
% Selección de líneas y cambios mecánicos explícitos; ninguna prueba nueva.
\subsection{Angular coordinates, turn normalization and action units}
\label{v:v:ant:sec:ii-angulos}

The action quantity and angular coordinate belong to different spaces. The section \(\hbar_{\rm ret}\) occupies the action line \(\mathcal L_S\); \(\eta_{\rm ret}\) is its coefficient in the declared basis. Angular expression acts on the dimensionless pair \((\alpha,\eta_{\rm ret})\) after its construction:
\begin{equation}
 \mathcal P(\alpha,\eta_{\rm ret})
 =\bigl([1000\alpha]_{360},[2(\eta_{\rm ret}+\alpha)]_{360}\bigr).
 \label{v:v:ant:eq:ii-par-angular}
\end{equation}
Nonadic completion determines the factor 1000; bidirectional closure determines the factor two. The circular chart of 360 units makes visible the partitions \(12\cdot30=4\cdot90=3\cdot120\). On the article's positive branch, the representatives are
\[
 A_{\deg}=1000\alpha,\qquad
 C^*_{\deg}=2(\eta_{\rm ret}+\alpha).
\]

\begin{proposition}[Shift of the base-thousand expression]
\label{v:v:ant:prop:ii-desplazamiento-angular}
Let \(\alpha=\sum_{n\geq1}b_n1000^{-n}\), with
\(b_n\in\{0,\ldots,999\}\), be the compatible expansion already generated.
The lifted coordinate \(A_{\deg}=1000\alpha\) satisfies
\[
 A_{\deg}=b_1+\sum_{n\geq1}b_{n+1}1000^{-n}.
\]
The expression preserves the sequence of blocks from the second onward; the initial block passes to the integer part.
\end{proposition}
\begin{proof}
The series converges absolutely. Multiplying every term by 1000 and separating the initial term gives the identity. The operation acts on the constructed expansion, with its blocks and memory; subsequent circular expression takes its class modulo 360.
\end{proof}

\begin{definition}[Charts of the angular coordinate]
The degree chart expresses a complete turn by 360 units. The radian chart expresses it by \(2\pi\). The coordinate changes and their inverses are
\[
 x=\frac{\pi}{180}A_{\deg},\qquad
 y=\frac{\pi}{180}C^*_{\deg},\qquad
 A_{\deg}=\frac{180}{\pi}x,\qquad
 C^*_{\deg}=\frac{180}{\pi}y.
\]
Along a path with preserved turn count, lifted coordinates are used; the circular quotient expresses its class modulo the corresponding turn.
\end{definition}

The definition characterizes degrees and radians by their exact relationship to the turn. The register also retains orientation, turns and transport memory. Expression of its angular class preserves phase, whereas the lift distinguishes paths having the same final phase. The physical reading of action must specify which of these coordinates enters the integral.

\begin{proposition}[Covariance of the action reader]
For \(h_a=2\pi\hbar_a\), the action of a lifted coordinate
satisfies
\begin{equation}
 \mathcal S_a(\theta)
 =\hbar_a\theta_{\rm rad}
 =h_a\frac{\theta_{\deg}}{360}.
 \label{v:v:ant:eq:ii-accion-cartas}
\end{equation}
A positive turn expresses \(h_a\).
\end{proposition}
\begin{proof}
Substituting \(\theta_{\rm rad}=\pi\theta_{\deg}/180\) gives \(\hbar_a\pi\theta_{\deg}/180
=(2\pi\hbar_a)\theta_{\deg}/360\). Over one turn \(\theta_{\deg}=360\), whence \(h_a\).
\end{proof}

The periodicity of a representation is declared separately. If the transported representation satisfies \(U(2\pi)=-I\), its operatorial restoration occurs at \(4\pi\). This property of the lift coexists with the per-turn normalization of \eqref{v:v:ant:eq:ii-accion-cartas}. Each preserves its object: integrated quantity, circular phase and represented state.

\subsection{Covariance of the channels and the incidence functional}

The channels are equivalently written as
\begin{equation}
 q_\pm=\exp[-(x\pm y)]
 =\exp\!\left[-\frac{\pi}{180}(A_{\deg}\pm C^*_{\deg})\right].
\label{v:v:ant:eq:ii-canales}
\end{equation}
\begin{proposition}[Recovery of the even and odd coordinates]
\label{v:v:ant:prop:ii-inversion-canales-angulares} In the lifted real chart, the positive channels uniquely determine the two coordinates:
\[
 A_{\deg}=-\frac{90}{\pi}\log(q_+q_-),\qquad
 C^*_{\deg}=\frac{90}{\pi}\log\frac{q_-}{q_+}.
\]
The exchange \((q_+,q_-)\mapsto(q_-,q_+)\) preserves \(A_{\deg}\) and reverses \(C^*_{\deg}\).
\end{proposition}
\begin{proof}
The sum and difference of the logarithms of \eqref{v:v:ant:eq:ii-canales} are, respectively, \(-\pi A_{\deg}/90\) and \(-\pi C^*_{\deg}/90\). The real logarithm is single-valued on the positive channels. The product is symmetric and the quotient is inverted when they are exchanged.
\end{proof}
In the twelve-position register, the antipodal involution \(r\mapsto r+6\pmod {12}\) forms six pairs. The uniform distribution of the odd coordinate over these pairs is \(C^*_{\deg}/6\): summing its six contributions recovers \(C^*_{\deg}\). This sectorial coordinate is distinct from the Planck constant \(h\), which belongs to the action line.

The conversion is applied only once to each argument. The quantity \(\alpha\) remains dimensionless and the quantity \(\hbar\) still belongs to the action line; their prior coordinates enter the map \(\mathcal P\).

The Barbero--Immirzi functional contains, in addition to the channels, a component bilinear in the angular representatives. Its transport is
\begin{equation}
 \frac{\pi}{180}A_{\deg}C^*_{\deg}
 =\frac{180}{\pi}xy.
 \label{v:v:ant:eq:ii-barbero-cartas}
\end{equation}
The identity is obtained by substituting both representatives in degrees with their expressions in radians. Thus, the complete formula in analytic coordinates is
\[
 \gamma_{\HMT}
 =\frac{180}{\pi}xy
 +12\bigl[S_{90}(e^{-x+y})-S_{90}(e^{-x-y})\bigr]
 -\bigl[S_{120}(e^{-x+y})-S_{120}(e^{-x-y})\bigr].
\]
Equivalence of the two charts transports the bilinear component and exponential functions simultaneously. This transformation specifies the precise role of the degree in the construction and makes evaluation reproducible without ambiguity of angular units.
